\section{The continuum problem}
\label{sec:continuum}

Part~\ref{part:fixed} constructs an infinite-volume theory with a
positive gap at fixed lattice spacing. This part is about the distance
between that and a continuum quantum Yang--Mills theory with a mass gap,
and it is written to be useful to someone trying to close that distance.

\subsection{The target, and what has no argument here}

The statement in question is the existence of a non-trivial quantum
field theory on $\R^{1,3}$ for $G = \SU(N)$, $N \ge 2$, satisfying the
Wightman axioms, with
\begin{equation}
  \spec(H)\big|_{\Omega^{\perp}} \subset [m, \infty),
  \qquad m > 0 .
\end{equation}

Four of its ingredients have no argument anywhere in this program and
are stated here so that no reader has to infer their absence.

\begin{enumerate}
  \item \textbf{Poincar\'e covariance.} No construction of a unitary
    representation of the covering group of $\mathcal{P}_+^{\uparrow}$ is
    given.
  \item \textbf{Restoration of rotational invariance.} The hypercubic
    lattice breaks it; nothing here restores it in the limit.
  \item \textbf{The continuum measure.} It is never constructed.
  \item \textbf{Reconstruction.} Even granting a limiting kernel, the
    passage to a Wightman theory is not carried out.
\end{enumerate}

What \emph{is} carried out is a set of estimates on the lattice side,
and this part states them, states why the six obstructions of
\S\ref{sec:breaks} block their assembly, and states in
\S\ref{sec:obligations} what would have to be proved.

\subsection{The three routes that are live}

The estimates in this program group into three approaches, which are
alternative rather than sequential. The ledger records them as
independent routes and explicitly notes that no chain of necessary
conditions runs between them.

\emph{The conditional-score route.} Control the conditional law of a
link given its complement, in the small-coupling limit, by dominating
the score of the conditional measure. The available inputs are strong:
exact heat-bath conditional laws, synchronized radial dilation following
the conditional minimum curve, the antipodal magnetic normal geometry,
exponential suppression outside the complete minimizing set, and a
uniform Agmon action gap on the complement. The obligation is
Problem~\ref{prob:m10}, and what stands between the inputs and the
conclusion is polynomial amplitude control and the tube moments, not a
missing structural idea.

\emph{The transport route.} Compare the fine and coarse theories by
transporting the vacuum and the sources along a common energy domain
against a physical clock. Here the recorded failure is instructive: an
earlier closure used a vacuum-only skew generator, and the source
generator transports the ground state rather than annihilating it. The
corrected obligation is Problem~\ref{prob:r10}, and its companion,
Problem~\ref{prob:grid}, is the requirement that the comparison's
constants not degrade with the number of interacting blocks.

\emph{The moving-time route.} Extract a limiting gap from approximate
sources and one growing observation time per cutoff. The implication
from its hypotheses to a gap is proved, and sharp: pairwise inequalities
give the optimal normalized rate, the horizon penalty and a witnessing
time, and one-atom examples show the rate cannot be improved. Its
hypotheses are the difficulty --- see Break~\ref{sec:break3}, which shows
the theorem returns exactly zero when the background error rate fails to
exceed twice the probe tilt rate, an inequality no calculation certifies
for four-dimensional Yang--Mills.

\subsection{Why none of the three currently closes}

Each route reaches a point where a summability or uniformity statement
is required, and at each such point the natural estimate is one power
short.

That is not a coincidence of technique. Break~\ref{sec:break1} exhibits
the arithmetic: on an asymptotically free trajectory the running
coupling gives increments of size $O(g_k^2) = O(1/k)$, the normalization
of the scale derivative carries no small factor, and $\sum 1/k$
diverges. Break~\ref{sec:break2} exhibits the same shortfall in the
polymer variable, where a contracting activity and a scale-independent
physical mass cannot be read off one schedule. Problem~\ref{prob:grid}
is the same shortfall again in the block count.

Underneath all three sits Break~\ref{sec:break5}: the estimates are
proved near $u = 0$ and the continuum lives at $u \to \infty$. A
summable increment would connect scales \emph{within} a regime; it would
still not move the construction between regimes. That is why
Problem~\ref{prob:trajectory} is marked load-bearing and the others are
not.
